\section{Motivation}\label{sec:motivation}

Large corpora are indispensable in contemporary political science, but the objects we ultimately analyze are not the raw texts themselves. Researchers want structured signals: whether an article reports a protest event; which actor sponsored an amendment; the date and location of a strike; or the count of distinct demonstrations in a monthly newsletter. We formalize this by positing an \emph{oracle} $\fstar$ that maps a full document $x$ to the task value $\fstar(x)\in\Y$ that the researcher actually cares about. In this view, LLMs are useful only insofar as they preserve---or at least accurately predict---the information contained in $\fstar(x)$; summaries are valid research objects only if they leave the oracle untouched. We therefore treat the summarizer $g$ as a targeted compression operator: its sole job is to discard tokens that are provably irrelevant to $\fstar$ while preserving every bit of evidence needed for the oracle to reach the same verdict.\footnote{Put differently, $g(x)$ must be a sufficient statistic for $\fstar(x)$. It may be terse or awkward, but it cannot omit the features that $\fstar$ inspects.} The oracle might be human annotation, a high-quality-but-expensive model, or a hypothetical infinite-context extractor;\footnote{In practice we approximate $\fstar$ with whatever expensive-but-accessible supervision we can afford—expert coders, deliberation-grade models, or hybrid workflows. The audit is therefore relative to that chosen oracle.} it is often too costly or infeasible to apply $\fstar$ to every raw document, which is precisely why we rely on compressed representations $g(x)$ in the first place.

This sufficiency framing is intentionally strong. When a task is non-local, $\fstar(b)$ for a single block $b$ may only be defined if the summary encodes the latent state required to process every possible continuation. We do not assume this decomposability for free; the audits in Section~\ref{sec:audit} and the diagnostics in Section~\ref{sec:examples} are precisely the procedures we use to test it. If those probes repeatedly fail, the remedy is to enrich the summary schema or enlarge the blocks until the oracle truly behaves like a sufficient statistic on substrings. Highlighting (and measuring) this modeling risk is therefore a first-class part of the framework, not a footnote.

Two practical constraints make this nontrivial. First, many documents are longer than a model's context window. In practice, one partitions $x$ into contiguous blocks, summarizes each block with an LLM, and recursively merges those summaries up a tree. Second, even when context is not the binding constraint, researchers often \emph{choose} to summarize for cost, latency, or workflow reasons (e.g., to store compact ``event summaries'' and relabel them as coding rules evolve). Either way, hierarchical summarization changes the object of analysis. The core methodological question is therefore: \emph{How can we guarantee this pipeline leaves $\fstar(x)$ unchanged without running the expensive oracle on the full text?}

Our answer is constructive as well as diagnostic. The structural invariants we introduce later double as optimization objectives: the same measurements that certify a pipeline also describe how to tune it. Frameworks such as DSPy \citep{khattab2024dspy} let us treat the summarization policy as a programmable system whose prompt parameters are updated until violations of these conditions disappear. Because each check inspects only the immediate child summaries, we obtain cheap supervision signals without new external labels.

\paragraph{Contributions.} Four claims structure the paper. First, we reinterpret mergeable sketches \citep{agarwal2013mergeable} as semantic invariance conditions on text, with Appendix~\ref{app:mergeable} showing that our sufficiency/idempotence/merge-consistency conditions reduce to mergeability in the commutative case. Second, Section~\ref{sec:mechanism} describes the algebraic ``ideal world'' in which summaries behave like homomorphisms on strings, giving us the normative target for pipeline design. Third, we introduce these three conditions (leaf-level sufficiency, on-range idempotence, and merge consistency) and treat them as prompt-engineering objectives that force a stochastic summarizer to act like a mergeable operator. Finally, Section~\ref{sec:audit} develops the probabilistic audit: a random search over the realized summary tree whose empirical violation rates drive both certification and the optimization loop described later. These ingredients yield an end-to-end blueprint for designing, validating, and tuning hierarchical summarizers.

Our focus is on \emph{decomposable} information-centric tasks such as event extraction, coding tuples, structured aggregation, or guided abstractive summaries with explicit rubrics where the oracle space $\Y$ is well defined on small spans. Open-ended creative summarization or stylistic rewriting---where $f^*$ rewards prose quality directly---lies outside our scope; in those settings the oracle metric itself is subjective, and our guarantees do not apply until the analyst turns the brief into a concrete oracle.
\section{Framework Overview}\label{sec:overview}

\subsection{From Global Algebra to Local Sampling}

Our answer relies on a shift in perspective: instead of trying to validate the final summary (which requires processing the full document to establish ground truth), we treat the summarization pipeline as a stochastic process and audit its atomic operations.

We introduce three consistency conditions: \emph{leaf sufficiency} (the summary preserves the oracle on small blocks), \emph{idempotence} (re-summarizing a summary is inert), and \emph{merge consistency} (pre-summarizing a span is equivalent to summarizing it jointly). The conditions can be satisfied by prompt design, DSPy-style inference-time optimization \citep{khattab2024dspy}, or fine-tuning; the details of their structural implications live in Appendix~\ref{app:proofs-core}. Appendix~\ref{sec:examples} instantiates these definitions for practical tasks including binary event detection, entity extraction, and numeric counting, illustrating how to audit each condition in practice.

This approach can be viewed as a semantic generalization of \emph{mergeable summaries} \cite{agarwal2013mergeable}. While that literature constructs algebraic data structures (sketches) to preserve commutative statistics (like counts or heavy hitters) under merge operations, we seek to preserve arbitrary semantic oracles $\fstar$ (like event extraction or narrative flow) under non-commutative string concatenation. We replace the explicit algebra of sketches with ``learned'' local consistency conditions that force a generic summarizer to behave \emph{as if} it were a mergeable sketch; Appendix~\ref{app:mergeable} shows that when $\fstar$ is symmetric, our conditions collapse exactly to the standard mergeability requirement.

The power of this framework lies in its efficiency. Because we prove that local consistency implies global preservation (Theorem~\ref{thm:one-pass}), we do not need to audit the root of the tree to monitor quality. Instead, we can perform a \emph{probabilistic audit}: by sampling and verifying a fixed number of leaves and internal merges---operations that fit easily within a small context window---we estimate the \emph{process error rates} that govern the hierarchy. This ``random search over the summary tree'' is auditable: every checked node runs entirely within context, yet the resulting failure-rate estimates expose how often the pipeline corrupts task values. The probabilistic audit therefore decouples the cost of quality control from document length when the goal is to monitor pipeline reliability, even though certifying the correctness of a specific book still requires extremely low local error rates.

\subsection{Optimizing for Consistency}

Beyond auditing, these conditions provide a recipe for optimization. Since the invariants are local, they can serve as objective functions for prompt engineering or fine-tuning (e.g., via frameworks like DSPy \citep{khattab2024dspy}). Rather than optimizing a model to ``summarize well''---a vague objective---practitioners can optimize the model to satisfy merge consistency: ``ensure that $g(\text{summary}_A + \text{summary}_B)$ yields the same extracted events as $g(\text{text}_A + \text{text}_B)$.'' Because each constraint operates on a single edge of the tree, gradients or prompt-adjustment heuristics can be estimated via the same random sampling procedure used for audits.

\subsection{Safe Learning on Summaries}

Finally, this local view yields a clean connection to learning. We consider supervision that depends on documents only through their task value, a common factorization that covers proper scoring, coding tasks, and Bradley--Terry--Luce models for pairwise preferences. Under this assumption, we show that training on oracle-preserving summaries is information-equivalent to training on full documents.

Specializing to Direct Preference Optimization (DPO) \citep{RafailovEtAl2023DPO}, we prove that if the hierarchy passes our audit, training a policy on the summaries achieves the same population optimum as training on the originals. When preservation is only approximate, the gap in the DPO loss is quantitatively controlled by the error rates estimated in our audit. This connects the mechanics of scalable preprocessing to the objectives of alignment: if the pipeline respects the three conditions, the preferences learned from the output are faithful to the source oracle, even when those preferences encode arbitrary, expensive judgments. Appendix~\ref{sec:dpo} provides the full derivation and extends the argument to any supervision scheme that factors through $\fstar$.
